\section{Related Work}

\subsection{Robotic Table Tennis}
Robotic table tennis has emerged as a rich testbed for high-speed perception, planning, and control, requiring rapid reactivity to return balls at competitive speeds.  In our own work, we were very much inspired by the  biologically-inspired approach to motion generation in \cite{mulling2010biomimetic, mulling2011biomimetic}. In this work the research team led by Jan Peters fits a template trajectory while minimizing the robot's deviations from a fixed ``comfort posture": Human demonstration via kinesthetic teach-in was used in \cite{mulling2013generalizestrike} to learn  a general mixture of motions policy from motion primitives. 
%A robot arm used model-free Reinforcement Learning to learn to smash balls without any real-world racket-ball contact during training \cite{buchler2022learning}. 
The interest in responding to adversaries  with  opponent modeling was presented in \cite{Wang_Boularias_Mülling_Peters_2011} with  anticipatory action selection \cite{6094892}, and  the design of the Intention-driven Dynamics Model \cite{Wang2013} and \cite{WANG2017399}. The work in \cite{MullingBMSP2014} applied inverse reinforcement learning to identify strategic elements of play from demonstrations between different players. Our own work has analyzed detailed play from hours of recorded video \cite{etaat2025latte} to determine intent. In  Section \ref{disc-sec}, we briefly discuss integrating this into our humanoid going forward. 

Recent progress has combined advances in hardware, control, and learning. For instance, a lightweight, high-torque robotic arm with model predictive control was developed to execute diverse hit styles with precision \cite{nguyen2025high}. Similarly, a model-based approach for hitting velocity control has demonstrated high success rates in returning balls accurately \cite{ji2021model}. On the learning side, seminal work on the use of fast learning-based techniques for playing table tennis have been presented by the Gemini Robotics team at DeepMind \cite{Sanketi2023, amor2025sas}. In continuation of this work, amateur human-level competitive performance was achieved through a hierarchical architecture with sim-to-real adaptation \cite{d2024achieving}. In contrast to large-scale training efforts, \cite{tebbe2021sample} showed that  reinforcement learning can learning to rally from scratch in under 200 trials. Moreover, beyond specialized robotic table tennis devices, general-purpose humanoid robots have been shown to play table tennis using impedance control \cite{xiong2012impedance}. However, in prior work, the humanoid was constrained to stand still without agile locomotion, which limited the effective hitting range. In our work, we focus on human-like and agile reaching motions for table tennis hitting, a capability not previously demonstrated on humanoid robots.

\subsection{Humanoid Whole-Body Control}
Whole-body humanoid control has seen rapid advances in recent years.
% single motion tracking
Early efforts mainly focused on training humanoid robots to mimic specific human motions using large-scale reinforcement learning \cite{cheng2024expressive, he2024omnih2o}. Other approaches decouple control into separate upper- and lower-body policies to simplify learning and coordination \cite{zhang2025falcon}. 
% general motion tracking
More recent research has shifted toward developing general whole-body motion trackers capable of imitating a variety of human movements \cite{ze2025twist, chen2025gmt}.
% ours
Building upon these advances \cite{liao2025beyond}, our work introduces a humanoid whole-body controller specifically designed for rapid and dynamic interactions with the environment, such as ball hitting.